%%%%%%%%%%  scan idx 446  |  folio 433  %%%%%%%%%%
complexes de Koszul associés. Le complexe de Koszul associé à la suite \(f_1,\ldots,f_p,g_1,\ldots,g_r\) est isomorphe au complexe produit tensoriel \(K_\bullet(u)\otimes_{\cO_X}K_\bullet(v)\). Comme c'est une résolution de \(\cO_Y\) par des \(\cO_X\)-Modules libres de type fini, les \(\underline{\Tor}_i^{\cO_X}(\cO_Y,\cO_Z)\) sont les faisceaux d'homologie du complexe \(K_\bullet(u)\otimes_{\cO_X}K_\bullet(v)\otimes_{\cO_X}\cO_Z\). On peut le considérer comme le produit tensoriel des complexes \(K_\bullet(u)\) et \(K_\bullet(v)\otimes_{\cO_X}\cO_Z\), et utiliser pour calculer son homologie les suites spectrales associées à un bicomplexe (EGA \(0_{\mathrm{III}}\) (11.3.5)), qui ont comme aboutissement l'homologie du complexe simple associé.

On considère la suite spectrale telle que \({}'E^2_{pq}=H^I_p\bigl(H^{II}_q(K_{\bullet\bullet})\bigr)\), avec les notations de EGA \(0_{\mathrm{III}}\) (11.3.5), le premier indice étant relatif au complexe \(K_\bullet(u)\) et le second au complexe \(K_\bullet(v)\otimes_{\cO_X}\cO_Z\). Or le complexe \(K_\bullet(v)\otimes_{\cO_X}\cO_Z\) est le complexe de Koszul sur \(\cO_Z\) associé à la suite de sections \(\overline g_1,\ldots,\overline g_r\), où les \(\overline g_i\) sont les images des \(g_i\) dans \(\cO_Z\), et ce complexe est une résolution de \(\cO_T\) sur \(\cO_Z\) car les immersions \(Z\to X\), et \(T\to Z\to X\) sont régulières (cf. démonstration de 1.7 ii). Comme \(\bigwedge^p E\) est libre sur \(\cO_X\), le complexe \(\bigwedge^pE\otimes_{\cO_X}(K_\bullet(v)\otimes_{\cO_X}\cO_Z)\) est une résolution de \(\bigwedge^pE\otimes_{\cO_X}\cO_T\). Par conséquent la suite spectrale dégénère, et les \(\underline{\Tor}_i^{\cO_X}(\cO_Y,\cO_Z)\) sont les faisceaux d'homologie du complexe \(K_\bullet(u)\otimes_{\cO_X}\cO_T\). Or la différentielle de \(K_\bullet(u)\) s'exprime en termes des \(f_i\), qui annulent \(\cO_T\); par suite, la différentielle de \(K_\bullet(u)\otimes_{\cO_X}\cO_T\) est nulle, et on obtient donc les isomorphismes:
\begin{equation}\tag{2.5.2}
\underline{\Tor}_i^{\cO_X}(\cO_Y,\cO_Z)
\xrightarrow{\sim}
\cO_T\otimes_{\cO_X}\bigwedge^iE
\longrightarrow
\bigwedge^i(\cO_T)^p.
\end{equation}

En particulier, on trouve pour \(i=1\):
\[
\underline{\Tor}_1^{\cO_X}(\cO_Y,\cO_Z)\xrightarrow{\sim}(\cO_T)^p,
\]
d'où pour tout \(i\):
\begin{equation}\tag{2.5.3}
\underline{\Tor}_i^{\cO_X}(\cO_Y,\cO_Z)
\xrightarrow{\sim}
\bigwedge^i\underline{\Tor}_1^{\cO_X}(\cO_Y,\cO_Z).
\end{equation}

